% Section 5 of the revised paper: evaluation protocol (draft 2026-10-25).

\section{Evaluation protocol}\label{sec:protocol}

An earlier version of this work compared PromptShift-CRC with three baselines on single runs of hand-built streams. On re-examination we found that a short-window baseline matches it, that single runs hide the seed-to-seed spread, and that the benchmark streams used no model outputs. This section describes what we do instead. Every real-stream setting was fixed before its results were computed; the synthetic constants are inherited from the earlier version and were chosen while those generators were built (Section~\ref{sec:limits}).

\paragraph{Metrics.}
Marginal coverage ($1-$ mean of $\mathrm{err}_t$) is reported, but Theorem~\ref{thm:longrun} bounds its distance from $1-\alpha$ for every online method, so it cannot rank them. The metrics that can are the following.
\begin{itemize}[nosep]
  \item \emph{Post-shift violation}: the mean of $\mathrm{err}_t$ over the first $L$ steps after each known regime boundary ($L=150$ for MMLU, $50$ for TruthfulQA, $150$ for the synthetic streams). Lower is better and the target is $\alpha=0.1$.
  \item \emph{Maximum local violation}: the largest mean of $\mathrm{err}_t$ over any window of $120$ steps ($60$ for TruthfulQA).
  \item \emph{Mean threshold}: the average of $\min\{q_t,1\}$, a measure of sharpness; a lower threshold means fewer unnecessary interventions.
  \item \emph{Intervention rate}: the fraction of steps with $q_t>\tau$, with $\tau=0.5$ fixed in advance.
  \item \emph{Frontier gap}: for each stream, Adaptive CRC is run with five windows. This traces a curve of post-shift violation against mean threshold. A method's frontier gap is its post-shift violation minus the value of that curve, linearly interpolated at the method's own mean threshold. A negative gap means the method does better than Adaptive CRC at the same sharpness.
\end{itemize}

\paragraph{Streams.}
\emph{Synthetic}: five generators, each with three regime changes, run with $30$ seeds ($0$ to $29$). These are the generators of the earlier version, with a deterministic per-task seed (the earlier script derived the seed from Python's per-process string hash, so its table could not be reproduced).
\emph{Real model outputs}: we run Qwen2.5-1.5B-Instruct on CPU with one fixed prompt template for each task and record the probability the model gives to each answer option. For MMLU \citep{hendrycks2021measuring} we use 45 randomly chosen questions per subject (a fixed seed), $2565$ in all, and for TruthfulQA \citep{lin2021truthfulqa} all $817$ questions in the multiple-choice setting, scoring each answer choice by the summed log-probability of its text. The risk is $1-p(\text{correct})\in[0,1]$, and as a second risk we use the 0/1 error. Features are unit-norm all-MiniLM-L6-v2 sentence embeddings of the question text \citep{reimers2019sentence}.
For RealToxicityPrompts \citep{gehman2020realtoxicityprompts} we draw $4000$ of the $99{,}442$ prompts at random (a fixed seed), continue each with GPT-2 (one nucleus sample with $p=0.9$ and $20$ new tokens), and score the continuation alone with the \texttt{unitary/toxic-bert} toxicity output; the features are embeddings of the prompt text, and the risk is that toxicity probability. A replication uses Qwen2.5-0.5B continuations scored by \texttt{unitary/unbiased-toxic-roberta} on the same prompts. The \emph{orderings} use only the benchmark's own metadata or the question text, never the model's output: random orders of the four broad MMLU categories ($24$ permutations), the $24$ permutations of the four quartile bands of the prompt toxicity annotation (RealToxicityPrompts; the annotation describes the prompt, not the model's output), ascending and descending prompt toxicity and of the $57$ subjects ($20$ orders), random orders of the TruthfulQA categories ($20$), embedding-direction orders (first principal component, forward and reversed), question-length order, and $20$ random permutations as a control with no shift.

\paragraph{Pre-declared settings.}
All methods use $\alpha=0.1$, the unprojected update with $\gamma$ fixed per experiment ($0.025$ for the synthetic pilot; $0.02$ for the other synthetic streams, MMLU and RealToxicityPrompts; $0.03$ for TruthfulQA, which has $n=817$), the correction factor $1+m_t$, and a pool of $3W$ steps for every weighted rule. On the real streams the calibration prefix and window are $256$ with a recent window of $128$ (MMLU) and $100$ with $50$ (TruthfulQA). The kernel bandwidth is the median pairwise distance $h$ among the first prefix prompts of that stream, the gate centre is $1.1h$ and the gate slope is $2.5/h$. These rules were written down before any real-stream result was computed. The synthetic-stream constants are those of the earlier scripts and were not changed; the earlier version chose them while developing the generators, which we say plainly in Section~\ref{sec:limits}.

\paragraph{Uncertainty.}
For synthetic streams we report the mean over $30$ seeds with the standard deviation and compare methods by paired differences within seed (mean difference, standard error, and the number of seeds in which one method is lower). For the real streams, the unit of replication is the ordering. For a single ordered stream, standard errors of coverage use a moving-block bootstrap of the violation indicator with block length equal to the recent window.
